\subsection{Correctness and Completeness}
\label{app:completeness-proof}
We prove Theorem~\ref{thm:complete-synthesis} under the assumptions of
Section~\ref{sec:formal-scope}, with every input dimension having extent
at least two. The proof has three parts: every valid
assignment passes the checks; passing the checks gives a valid construction;
and the restricted additive fragment separates into independent components.

\subsubsection{Necessity}
\begin{lemma}[Necessity]
\label{lem:necessity}
For one component,
$\exists\beta:\operatorname{Valid}(m,\beta)\Rightarrow\operatorname{check}(m)$.
\end{lemma}
\begin{proof}
Suppose $\beta$ is valid. Lemma~\ref{lem:canonicalization} gives a valid
canonical assignment $\beta'$ with live set $S_v$ at every node. Its final
output agrees with the LHS, so
$S_{\mathrm{root}}=\operatorname{set}(F)$.

Starting with the root requirement $F$, prove from parents to children that
each canonical output contains $Req_v$ as a subsequence. A contraction's
output consists of batch, left-free, and right-free blocks in that order.
Any subsequence of that output must follow the same block order, giving
Equation~\eqref{eq:block-order}. Each free subsequence must occur in the
corresponding child. Batch order can be chosen independently of child order.

Unary summation preserves the order of surviving axes, so its child contains
the same required subsequence. A permutation imposes no child requirement.
These are exactly the propagation rules of
Appendix~\ref{app:operation-construction-rules}. At an input, the propagated
requirement must be a subsequence of its actual order, $Req_v\preceq I_v$.
Thus $C_v$ is true at every node. Together with the root live-set check,
this gives $\operatorname{check}(m)=\mathrm{true}$.
\end{proof}

\subsubsection{Sufficiency}
\begin{lemma}[Sufficiency]
\label{lem:sufficiency}
For one component,
$\operatorname{check}(m)\Rightarrow
\operatorname{Valid}(m,\operatorname{build}(m))$.
\end{lemma}
\begin{proof}
Assume the checks pass. By structural induction on the RHS, we maintain:
\begin{enumerate}
\item The constructed attributes and intermediate operations are legal.
\item The output axes are $S_v$ and satisfy $Req_v$.
\item The input-port partition equals the LHS partition restricted to
the subtree's inputs, including classes removed from the output.
\end{enumerate}
The second invariant is Equation~\eqref{eq:order-invariant}.
An input uses its fixed order, which passes the input subsequence check.
A permutation freely arranges its axes. Unary summation removes exactly
the specified axes and preserves the rest.

At a contraction, the construction uses $B_v,K_v,LF_v,RF_v$.
The block-order check ensures that placing required batch axes first and
inheriting the child free orders satisfies $Req_v$. The output axes are
exactly $S_v$. Each selected axis occurs once in its child, so all position
lookups are defined. Paired axes have equal extents because they belong to
the same LHS class. Batch and contraction sets are disjoint. Thus all
constructed attributes are legal.
By D4, the node assignments combine without conflicting choices.

For the third invariant, consider the input ports from $U_v$.
At an input, distinct axes belong to distinct classes. Permutation and
unary summation do not merge ports. At a contraction, any LHS class with
support in both child subtrees is mandatory in both children, so it appears
in both live sets. The contraction connects these two occurrences and never
connects distinct LHS classes. The combined partition is therefore exactly
the LHS partition restricted to the current subtree.

A class is removed only when it is absent from the final output and all its
support is already inside the subtree. The partition invariant ensures that
its supporting ports have been connected, so no later connection is lost.
At the root, the full input partition consequently agrees with the LHS.

Finally, $S_{\mathrm{root}}=\operatorname{set}(F)$ and
$F\preceq\operatorname{Ord}_{\mathrm{root}}$. The output has exactly these
axes, each once, so $\operatorname{Ord}_{\mathrm{root}}=F$. Hence
$N(L)=N(R)$. Lemma~\ref{lem:axis-exactness} gives equality for all input
values. Together with legality, this proves
$\operatorname{Valid}(m,\operatorname{build}(m))$.
\end{proof}

\subsubsection{Extension to Multiple Terms}
\label{sec:addition-proof}
We extend the single-component result to sums of components under
Section~\ref{sec:formal-scope}'s additive restrictions. Write
\begin{equation}
 L=\sum_{t\in J}L_t,\qquad R=\sum_{t\in J}R_t.
\end{equation}
Expanding the LHS preserves captured attributes and replaces operands with
summands of the same shape. It therefore preserves operation legality and
output shape.

To see why equality separates by component, multiply each input $X_i$ by
an independent real scalar $z_i$. Read-once components are multilinear, so
\begin{equation}
 L_t((z_iX_i)_i)=\left(\prod_{X_i\in U_t}z_i\right)L_t(X).
\end{equation}
The same identity holds for $R_t$. Fix tensor values and an output
coordinate, and compare the two sides as polynomials in the $z_i$.
Distinct input sets $U_t$ give distinct monomials. Uniqueness of polynomial
coefficients therefore gives, for legal expressions with equal output shape,
\begin{equation}
 \forall X,\ L(X)=R(X)
 \quad\iff\quad
 \forall t\in J,\ \forall X,\ L_t(X)=R_t(X).
 \label{eq:additive-separation}
\end{equation}
Apply the component algorithm independently and combine the results:
\begin{equation}
 \operatorname{check}(m)=\bigwedge_{t\in J}\operatorname{check}_t(m),
 \qquad
 \operatorname{build}(m)=\bigcup_{t\in J}\operatorname{build}_t(m).
 \label{eq:additive-construction}
\end{equation}
The RHS components use disjoint attribute identifiers, so the assignments
cannot conflict. Each successful construction has its LHS component's
output shape. LHS expansion preserves shape, so all these shapes agree
with the overall LHS, making the RHS additions legal.

Each component has its own axis partition. Inputs may appear in several
terms, but addition only aligns output coordinates; it does not merge input
port classes across terms.

The restrictions are necessary for this argument. Terms with the same
input set cannot be separated by the scaling variables. D4 ensures that
component assignments combine without conflict; its boundary example
appears in Table~\ref{tab:direct-scope}.

Lemmas~\ref{lem:necessity} and~\ref{lem:sufficiency} prove the theorem for
one component. Equation~\eqref{eq:additive-separation} and the independent
assignment construction extend both directions to the additive fragment.
This completes the proof of Theorem~\ref{thm:complete-synthesis}.
